\section{Distance Estimation}\label{distance-estimation}\index{distance estimation}

Distance Estimation (DE) is the calculation of an estimated distance from the
given point to the nearest surface of the fractal. As suggested by the word
'estimate', it is an approximate value. It is calculated using simplified
algorithms based on analytical (Analytical DE)\index{distance estimation!analytical} or numerical (Delta DE)\index{distance estimation!delta DE}
calculations of gradients.

DE is the most important algorithm required to render three-dimensional fractals
within a reasonable time.  A ray is traced from the camera for each pixel (1000 x 1000 resolution = 1,000,000 rays). They match FOV from the camera eye (i.e., they are not
parallel).

The DE provides an approximation of where along a ray the fractal surface should be located. This provides a starting point for stepping along the ray while looking for the surface location.
The process of stepping along the ray and testing for the surface location, is called ray-marching.

Without the DE calculation, the proximity of the fractal surface would need to be repeatedly calculated after each of many very small steps along the ray from the camera. For example, without an estimate of where the fractal surface is, you may need up to 10,000 steps to trace a ray of light, for every pixel of the image.

Ray marching\index{ray marching} looks like the illustration below. In each step, an estimation on
the distance to the nearest fractal surface is calculated. The calculation point is moved
along the ray by this distance. The next step is re-calculated based on the
estimated distance. This distance is smaller so this time the step is a
smaller distance. The ray-marching becomes more accurate closer to the surface
of the fractal. The ray marching ceases when a step moves the calculation point to be within a set
``distance threshold''\index{ray marching!distance threshold} from the surface or after a maximum number of iterations
if the option ``stop at maximum iteration'' is enabled.

Since the estimation contains some error (sometimes quite large), there is a risk that a step will be too large, and incorrectly step past the fractal surface. This may result in visible noise in the rendered image.

To prevent this, the step can be multiplied by a number between 0 and 1 (\emph{ray-marching step multiplier}\index{ray marching!step multiplier}). Steps are then smaller, so there is less risk of ``overstepping'' the surface (better image quality), but the rendering time increases due to more steps being required (figure \ref{distance_estimation_defactor_1}, \ref{distance_estimation_defactor_05}).

\simpleImageWithCaptionHalfWidth{img/manual/media/distance_estimation_defactor_1.png}
{Distance Estimation with DE factor 1}
{distance_estimation_defactor_1}{h}

\simpleImageWithCaptionHalfWidth{img/manual/media/distance_estimation_defactor_05.png}
{Distance Estimation with DE factor 0.5}
{distance_estimation_defactor_05}{h}

\subsection{Analytical DE vs Delta DE}\label{analytical-vs-delta-de}

Each formula is assigned a DE type: either \emph{Analytical DE} or \emph{Delta DE}.

\emph{Analytical DE}\index{distance estimation!analytical} computes the derivative of the fractal formula mathematically. It tracks the derivative of the iteration function alongside the point position, producing an accurate distance estimate. This mode is faster because the derivative is computed directly during iteration.

\emph{Delta DE}\index{distance estimation!delta DE} approximates the derivative numerically by comparing the fractal function evaluated at two nearby points. This method does not require a hand-crafted derivative formula, making it more robust for formulas with complex or discontinuous transformations. However, it is slower because each evaluation requires computing the fractal function twice.

In most cases, Analytical DE produces good results and is the preferred mode for speed. However, with some formulas (particularly those with discontinuous folds or conditional transforms), only Delta DE will produce a correct image. If you observe artifacts or noise that cannot be fixed by adjusting the step multiplier, try switching to Delta DE.

\subsection{DE Functions}\label{de-functions}

When using DE (either Analytical or Delta), the final distance is calculated by combining the radius \emph{r} (the distance of the final iteration point from the origin) with the computed derivative \emph{DE}. The formula used depends on the DE function type, which is determined by the nature of the fractal's transformations:

\begin{description}
	\item[Linear DE] Used for formulas with linear transformations (e.g., Mandelbulb, Mandelbox). The distance is calculated as:
	\[ distance = \frac{r}{\lvert DE \rvert} \]
	where \emph{r} is the radius (distance from the origin to the final iteration point) and \emph{DE} is the computed derivative.
	
	\item[Logarithmic DE] Used for formulas with logarithmic transformations. The distance is calculated as:
	\[ distance = \frac{0.5 \cdot r \cdot \log(r)}{DE} \]
	This function is applied when \emph{r} > 1.0. For \emph{r} $\le$ 1.0, the distance is set to 0.0.
	
	\item[IFS DE] Used for IFS (Iterated Function System) formulas such as Menger Sponge and Sierpinski. The distance is calculated as:
	\[ distance = \frac{r - 2.0}{\lvert DE \rvert} \]
	The constant 2.0 accounts for the specific scaling properties of IFS fractals.
	
	\item[Pseudo-Kleinian DE] Used for Pseudo-Kleinian fractals. The distance is calculated as:
	\[ distance = \frac{\max(r_{xy} - DE_{pseudoKleinian}, \lvert r_{xy} \cdot z.z \rvert / r)}{\lvert DE \rvert} \]
	where $r_{xy} = \sqrt{z.x^2 + z.y^2}$.
	
	\item[Jos-Kleinian DE] Used for Jos Kleinian fractals. The distance is calculated using a modified formula that accounts for the sphere folding parameters.
	
	\item[Max Axis DE] An alternative DE function that uses the maximum absolute component of the iteration point:
	\[ distance = \frac{\max(\lvert z.x \rvert, \lvert z.y \rvert, \lvert z.z \rvert)}{\lvert DE \rvert} \]
	
	\item[Custom DE] Some formulas provide a custom distance calculation via the \emph{aux.dist} auxiliary variable.
\end{description}

The automatic DE function is assigned based on the formula type. These settings can be overridden manually on the \emph{Render Engine} tab if needed.

\subsection{Linear DE offset}\label{linear-de-offset}

For hybrid fractals that mix mandelbox-type and IFS-type formulas, the \emph{Linear DE offset} parameter adjusts the DE calculation. It is located in the Hybrids tab.

The offset modifies the linear DE formula as follows:
\[ out.distance = \frac{r - offset}{\lvert DE \rvert} \]

The Linear DE offset is normally used in the range from 0.0 (mandelbox) to 2.0 (IFS). Adjusting this parameter can assist in fine-tuning the DE calculation when mixing these formula types in a hybrid fractal.

\simpleImageWithCaptionHalfWidth{img/manual/media/linear_de_offset.png}
{Linear DE offset parameter}
{linear_de_offset}{h}

\subsection{Bad DE statistic}\label{bad-de-statistic}

\simpleImageWithCaptionHalfWidth{img/manual/media/dock_statistics.png}
{Statistics Tab with histogram data}
{dock_statistic}{h}

In the Statistics (figure \ref{dock_statistic}) (enable in \emph{View} menu) you can see \emph{Percentage of
	Wrong Distance Estimations} (``Bad DE''). The statistic is only available in no-OpenCL mode. This number is the percentage of image
pixels which potentially have big errors in distance estimation calculation
(estimated distance was much too high). It is visible as a noise on the image.
As a general rule less than 0.01 is good, but it is case specific and 3.0
sometimes is OK and 0.0001 sometimes is not.

Figure \ref{bad_de_menger_sponge} is an example of Bad DE due to over-stepping. It starts to occur at the corner nearest to the camera, resulting in the black areas and disintegration of the fractal. If the ray-marching step multiplier is set even higher, the fractal will entirely disintegrate. This disintegration will generally be reflected in the Percentage of Wrong Distance Estimations statistic.

\simpleImageWithCaptionHalfWidth{img/manual/media/bad_de_menger_sponge.png}
{Example of over-stepping}
{bad_de_menger_sponge}{h}

\subsection{Troubleshooting DE issues}\label{de-troubleshooting}

If you observe artifacts or noise in your rendering, try the following:

\begin{enumerate}
	\item \textbf{Increase the ray-marching step multiplier} (decrease the value) to reduce the risk of overstepping the surface.
	\item \textbf{Switch DE mode} from Analytical to Delta (or vice versa) if the current mode produces incorrect results for the given formula.
	\item \textbf{Adjust the Linear DE offset} when using hybrid fractals that mix mandelbox and IFS formulas.
	\item \textbf{Monitor the Bad DE statistic} in the Statistics dock to check if DE quality improves after changes.
\end{enumerate}

The quality produced by the DE mode and function combinations is formula-specific. The setting of formula parameters can also greatly affect the quality produced by the DE. In some cases the choice of fractal image is determined by
what location and parameters can produce good DE quality.
